% answers_short.tex —— 2010 年·普通高等学校招生考试·重庆卷·文科 · 简短答案速查
% 由 tools/make_short_answers.py 自 answers.tex 的 \gksolhead/\gkanswer 提取，
% 并与 2010/2010_数学真题_重庆文_答案要点.md 逐题比对后生成，勿手改（改源 MD 后重跑）。

\begin{gknotice}
\textbf{使用说明}：本节只给最终答案便于快速核对，完整推导见解析版。答案为编写者独立推导，\textbf{非官方标准答案}。
\end{gknotice}

\gkpart{选择题}
\begin{center}\begin{tabular}{ccccccccccc}
\toprule
题号 & 1 & 2 & 3 & 4 & 5 & 6 & 7 & 8 & 9 & 10 \\
\midrule
答案 & B & A & D & C & B & A & C & D & D & C \\
\bottomrule
\end{tabular}\end{center}

\gkpart{填空题}
\begin{itemize}[leftmargin=2.6em]
  \item[\textbf{11.}] $(-1,0)$。
  \item[\textbf{12.}] $-2$。
  \item[\textbf{13.}] $2$。
  \item[\textbf{14.}] $\dfrac{3}{70}$。
  \item[\textbf{15.}] $-\dfrac12$。
\end{itemize}

\gkpart{解答题}
\begin{itemize}[leftmargin=2.6em]
  \item[\textbf{16.}] （1）$a_n=21-2n$，$S_n=20n-n^2$；（2）$b_n=3^{n-1}+21-2n$，$T_n=\dfrac{3^n-1}{2}+20n-n^2$。
  \item[\textbf{17.}] （1）$\dfrac15$；（2）$\dfrac23$。
  \item[\textbf{18.}] （1）$\dfrac13$；（2）$-\dfrac72$。
  \item[\textbf{19.}] （1）$f(x)=-\dfrac13x^3+x^2$；（2）增区间 $(-\sqrt2,\sqrt2)$，减区间 $(-\infty,-\sqrt2)$，$(\sqrt2,+\infty)$；在 $[1,2]$ 上最大值 $\dfrac{4\sqrt2}{3}$，最小值 $\dfrac43$。
  \item[\textbf{20.}] （1）证明；（2）$-\dfrac{\sqrt3}{3}$。
  \item[\textbf{21.}] （1）$\dfrac{x^2}{4}-y^2=1$，渐近线 $y=\pm\dfrac12x$；（2）$3$。
\end{itemize}
